% !PRINTED 324
%<<CONT
be put in place of $\varphi$. The two arbitrary quantities $c$ and $d$
remain in it.

If, for the sake of a more elegant notation, we further set $c = -\eta$,
$d = \xi$, we thus obtain:

\weq{8}{S\left(\frac{\xi\beta - \eta\alpha}{x\beta - y\alpha}\right)^2 = \frac{[\xi f'(x) + \eta f'(y)]^2 - (x\eta - \xi y)^2 H(x,y)}{nf(x,y)^2},}

and in this form both the right-hand and the left-hand side remain
completely unchanged if we apply an arbitrary linear transformation
simultaneously to the pairs of variables $x, y$; $\xi, \eta$; $\alpha, \beta$.

The formula (8) is an identity; it holds for every function
$f(x, y)$ and for all values of the variables $x, y$; $\xi, \eta$. But now
we make the assumption that $f(x, y)$ has only real roots, that is,
that the $\alpha, \beta$ are real, and we also suppose $\xi, \eta$ and
$x, y$ to be real. Both sides of the equation (8) are then
essentially positive. We denote their common value by
$P$, and now determine a quantity $X:Y$ by the quadratic
equation
\[
P = \left(\frac{X\eta - Y\xi}{Xy - Yx}\right)^2,
\]
or

\weq{9}{\Phi = (Xy - Yx)^2 P - (X\eta - Y\xi)^2 = 0.}

This equation has two real roots, which are obtained from
\[
X\eta - Y\xi = \pm \sqrt{P}\,(Xy - Yx)
\]
About the position of the roots of this equation we can, however,
assert the following:

If we set $X = x$, $Y = y$, then $\Phi$ becomes negative.

But if we set $X = \alpha$, $Y = \beta$, where $\alpha:\beta$ is any one
of the roots of $f = 0$, then
\[
\left(\frac{X\eta - Y\xi}{Xy - Yx}\right)^2 = \left(\frac{\xi\beta - \eta\alpha}{x\beta - y\alpha}\right)^2,
\]
that is, equal to \gesperrt{one term} of the sum
\[
P = S\left(\frac{\xi\beta - \eta\alpha}{x\beta - y\alpha}\right)^2
\]
and therefore smaller than $P$. From this it follows that $\Phi$ \gesperrt{is
positive} for $X:Y = \alpha:\beta$.

If we now arrange the values of all the $\alpha:\beta$ and $x:y$,
ordered cyclically by magnitude, say on a circle,
%CONT>>
